\section*{第 \ref{ch:b1:body-plans-tissues} 章——动物的身体构型与组织}

\begin{solution}{exo:b1:body-plans-tissues:1}
\omterm{def:b1:body-plans-tissues:symmetry}{辐射对称}：围绕一个轴排列，没有左右之分（水螅、
水母）。\omterm{def:b1:body-plans-tissues:symmetry}{两侧对称}：有一个镜面，区分前后与背腹
（蚯蚓、青蛙）。\omterm{def:b1:mammal-organization:bodyplan}{头部集中化}与\omterm{def:b1:body-plans-tissues:symmetry}{两侧对称}相伴：头先行移动的动物，
其前端迎接外界，感觉\omterm{def:b1:mammal-organization:organ}{器官}
与神经中枢在此集中。
\end{solution}

\begin{solution}{exo:b1:body-plans-tissues:2}
\omterm{def:b1:body-plans-tissues:layers}{外胚层}：皮肤\omterm{def:b1:flowering-plant-organization:tissues}{表皮}、\omterm{def:b1:body-plans-tissues:nervous}{神经组织}。\omterm{def:b1:body-plans-tissues:layers}{中胚层}：肌肉、骨、
血液、\omterm{def:b1:body-plans-tissues:connective}{结缔组织}。\omterm{def:b1:body-plans-tissues:layers}{内胚层}：消化道的内壁。
\end{solution}

\begin{solution}{exo:b1:body-plans-tissues:3}
\omterm{def:b1:body-plans-tissues:epithelium}{上皮}\omterm{def:b1:body-plans-tissues:tissue}{组织}：排列紧密的极性细胞构成片层，位于\omterm{def:b1:body-plans-tissues:epithelium}{基膜}之上，
起覆盖、吸收和分泌作用。\omterm{def:b1:body-plans-tissues:connective}{结缔组织}：细胞散布于丰富的基质中，
起连接、支持和运输作用。\omterm{def:b1:body-plans-tissues:muscle}{肌肉组织}：细胞内充满收缩丝，
起运动作用。\omterm{def:b1:body-plans-tissues:nervous}{神经组织}：由\omterm{def:b1:body-plans-tissues:nervous}{神经元}和\omterm{def:b1:body-plans-tissues:nervous}{神经胶质细胞}组成，
起感受、整合和指挥作用。
\end{solution}

\begin{solution}{exo:b1:body-plans-tissues:4}
紧密连接（封闭细胞间隙）、黏着连接
（与肌动蛋白环带相连，起固定作用）、桥粒（机械性铆钉，
与中间纤维相连）、缝隙连接（两侧细胞质之间的
通道）。
\end{solution}

\begin{solution}{exo:b1:body-plans-tissues:5}
三\omterm{def:b1:body-plans-tissues:layers}{胚层}、有\omterm{def:b1:body-plans-tissues:layers}{体腔}、具贯通的消化道、身体\omterm{def:b1:body-plans-tissues:segments}{分节}，
体壁柔软，依靠\omterm{def:b1:body-plans-tissues:segments}{流体静力骨骼}：环节动物（蚯蚓）。
\end{solution}

\begin{solution}{exo:b1:body-plans-tissues:6}
每一体节内的\omterm{def:b1:body-plans-tissues:layers}{体腔}液不可压缩。环肌收缩时，该节
变细变长（向前推进）；纵肌收缩时，该节变短变粗
（借刚毛锚定于洞壁）。两种收缩沿身体
向后传播的波，
将蚯蚓向前拉动。
\end{solution}

\begin{solution}{exo:b1:body-plans-tissues:7}
\omterm{def:b1:body-plans-tissues:segments}{外骨骼}：坚硬、起保护作用，肌肉附着于其内侧；不能生长，
因此动物必须蜕皮，蜕皮期间柔软易损；在陆地上，其质量随
体长的三次方增加，而强度仅随二次方增加，
这限制了节肢动物的体型只能较小。\omterm{def:b1:body-plans-tissues:segments}{内骨骼}：随动物一同生长，
无须蜕皮，表面保护较弱，肌肉附着于其外侧；
允许体型较大。
\end{solution}

\begin{solution}{exo:b1:body-plans-tissues:8}
紧密连接迫使一切物质都须穿过细胞，其细胞膜上
带有选择性转运蛋白；\omterm{def:b1:body-plans-tissues:epithelium}{上皮}由此能够选择何种物质通过、
以何种方向通过。若无紧密连接，管腔内容物会经细胞间隙
渗入\omterm{def:b1:organism-environment:milieu}{内环境}（细菌、毒素、未消化的分子），
已吸收的营养物质也会漏回。
\end{solution}

\begin{solution}{exo:b1:body-plans-tissues:9}
小肠。黏膜（带绒毛的\omterm{def:b1:body-plans-tissues:epithelium}{上皮}、其下的疏松\omterm{def:b1:body-plans-tissues:connective}{结缔组织}、
薄的黏膜肌层），黏膜下层（致密\omterm{def:b1:body-plans-tissues:connective}{结缔组织}），
肌层（内层环行、外层纵行的平滑肌），
浆膜。
\end{solution}

\begin{solution}{exo:b1:body-plans-tissues:10}
软骨细胞依靠凝胶中的扩散，从周围血管获取营养，
因此该\omterm{def:b1:body-plans-tissues:tissue}{组织}较薄或代谢缓慢，愈合能力差，
用于需要光滑、可压缩、无血管表面的部位，
如关节面、耳、鼻以及骨的生长端。骨细胞位于
不容物质扩散通过的矿化基质之中，须由管道中的血管供养；
骨愈合良好，终生持续重建，
构成承重的\omterm{def:b1:body-plans-tissues:segments}{骨骼}。
\end{solution}

\begin{solution}{exo:b1:body-plans-tissues:11}
吸收失去选择性：营养物质沿浓度梯度在细胞间通过，
转运蛋白的定向作用随之丧失，而已吸收的分子
漏回管腔。\omterm{def:b1:organism-environment:milieu}{内环境}接纳管腔中的一切内容，即
离子、水、未消化的多肽与细菌产物，同时水分和离子
渗入肠道：腹泻、水肿或脱水，\omterm{def:b1:mammal-organization:compartments}{血浆}成分紊乱。细菌与毒素
穿过屏障，
引起炎症和感染。
\end{solution}

\begin{solution}{exo:b1:body-plans-tissues:12}
构型是遗传而来的，并非为适应动物的需要而选择，
它决定了此后演化的可能方向：\omterm{def:b1:body-plans-tissues:segments}{外骨骼}要求蜕皮，
并使昆虫保持小体型；脊椎动物必须由体节构建身体，
即便并不需要它们；\omterm{def:b1:body-plans-tissues:symmetry}{辐射对称}的动物难以演化出头部。
表中列出的是各谱系所得到的备选方案，每一种在其约束之内
都完整而成功——它们并非阶梯，
青蛙的构型并不优于小龙虾，
只是可能性的另一组组合。
\end{solution}

\begin{solution}{pb:b1:body-plans-tissues:1}
\textbf{1.} 黏膜：\omterm{def:b1:body-plans-tissues:epithelium}{上皮}、疏松\omterm{def:b1:body-plans-tissues:connective}{结缔组织}、平滑肌薄层。
黏膜下层：致密\omterm{def:b1:body-plans-tissues:connective}{结缔组织}（含血管、神经）。
肌层：两层平滑肌，其间为\omterm{def:b1:body-plans-tissues:nervous}{神经组织}。
浆膜：\omterm{def:b1:body-plans-tissues:connective}{结缔组织}，表面覆以\omterm{def:b1:body-plans-tissues:epithelium}{上皮}。
\textbf{2.} 小肠；胃黏膜无绒毛，但具
深陷的管状腺。
\textbf{3.} 肠道自身神经网络的\omterm{def:b1:body-plans-tissues:nervous}{神经元}胞体；它们控制
收缩（混合与推进）和分泌。
\textbf{4.} 环肌层使管腔变窄（\omterm{def:b1:body-plans-tissues:segments}{分节}运动、挤压）；
纵肌层使其缩短。两者交替收缩，
推动内容物向前移动。
\textbf{5.} $\pi\times 0.03\times 6 = \qty{0.57}{m^2}$。
\textbf{6.} \qty{1.7}{m^2}。
\textbf{7.} $\pi\times 0.15\times 0.6 = \qty{0.28}{mm^2}$。
\textbf{8.} $25\times 0.28 = \qty{7.1}{mm^2}$ 即每 mm$^2$，倍数为 $7$。
\textbf{9.} $1.7\times 7.1 = \qty{12}{m^2}$。
\textbf{10.} $\pi\times 0.1\times 1 = \qty{0.31}{\micro m^2}$；
\qty{470}{\micro m^2} 每个细胞。
\textbf{11.} 顶端面 \qty{25}{\micro m^2}：倍数为 $19$。
\textbf{12.} $12\times 19 = \qty{230}{m^2}$。
\textbf{13.} $\qty{0.28}{mm^2}/\qty{25}{\micro m^2} = \num{11000}$ 个细胞。
\textbf{14.} $1.7\,\mathrm{m^2}\times 25\times 10^6\,\mathrm{mm^{-2}} =
\num{4.2e7}$ 个绒毛；$\num{4.7e11}$ 个细胞。
\textbf{15.} 每天 $\num{1.2e11}$ 个，每秒 $\num{1.4e6}$ 个。
\textbf{16.} 位于绒毛基部的腺体（隐窝）中；
分裂中的细胞受到底部的保护，其产物沿绒毛
向上迁移，途中吸收营养，并从顶端脱落。
\textbf{17.} $12\,\mathrm{m^2}\times 25\times 10^{-6}\,\mathrm{m} =
\qty{3e-4}{m^3} = \qty{300}{cm^3}$，约 \qty{315}{g}。
\textbf{18.} 周长 \qty{20}{\micro m}，即每个细胞 \qty{10}{\micro m}
（每条边由两个细胞共享）；$\num{4.7e11}\times 10^{-5} = \qty{4700}{km}$。
\textbf{19.} 紧密连接封闭了细胞旁路径；
物质必须经顶端膜进入细胞，
再经基底膜离开。
\textbf{20.} \omterm{def:b1:body-plans-tissues:connective}{胶原蛋白}和\omterm{def:b1:body-plans-tissues:connective}{弹性蛋白}，由成纤维细胞分泌。
\textbf{21.} \omterm{def:b1:body-plans-tissues:epithelium}{上皮}：\omterm{def:b1:body-plans-tissues:layers}{内胚层}；核心（\omterm{def:b1:body-plans-tissues:connective}{结缔组织}、血管、
肌肉）：\omterm{def:b1:body-plans-tissues:layers}{中胚层}。
\textbf{22.} 泵送作用：绒毛缩短时挤压淋巴管和
毛细血管，更新核心中的液体，使梯度
穿过\omterm{def:b1:body-plans-tissues:epithelium}{上皮}时保持陡峭。
\textbf{23.} 因绒毛倍数降至 $7$（扁平\omterm{def:b1:body-plans-tissues:epithelium}{上皮}的微绒毛
仍存在）：脂肪、糖类和维生素吸收不良，体重
下降，腹泻。
\textbf{24.} $0.28\times 19 = \qty{5.4}{mm^2}$。
\textbf{25.} 每个绒毛约有 \qty{5}{mm^2} 的吸收膜，
绒毛高 \qty{0.6}{mm}，小肠中约有
四千万个绒毛。
\end{solution}
